\section{Introduction}

Traffic-signal controllers are commonly trained and tested within one fixed
network.  New random seeds change arrivals, but topology, route choice, phase
programs and feature semantics remain largely unchanged.  Cross-network
deployment changes all of them at once.  A policy or world model can therefore
fit source transitions while misranking actions in a target network, which is
the quantity that ultimately determines control performance.

Two established approaches expose the central trade-off.  MaxPressure and
phase-pressure controllers use local queue imbalance and transfer without a
city-indexed policy, with strong throughput and stability motivation
\citep{varaiya2013maxpressure}.  Pressure structure also underlies learned
controllers including PressLight, FRAP and MPLight
\citep{wei2019presslight,zheng2019frap,chen2020thousandlights}.  These rules are
difficult baselines because local pressure is already close to the control
sufficient statistic in many signal networks.  Conversely, simulator-based
offline-to-online reinforcement learning can use source data to correct
dynamics mismatch \citep{niu2022h2o,liu2025h2oplus}, but an unrestricted
residual may absorb network identity, demand scale and phase-specific
associations that do not survive a city shift.

Prior transfer-oriented signal control spans materially different information
budgets.  MetaLight adapts a meta-learned initialization with target
experience, whereas GESA and X-Light construct scenario-normalized or
cross-city representations for policy generalization
\citep{zang2020metalight,jiang2024gesa,jiang2024xlight}.  CrossLight explicitly
combines source-city offline pretraining with target-city online fine-tuning
\citep{sun2024crosslight}.  Sim-to-real work such as UGAT instead transforms
actions to compensate for a simulator dynamics gap \citep{da2023ugat}.  These
methods establish that topology normalization, meta-adaptation and dynamics-gap
correction are each useful, but their zero-shot, online-interaction and
sim-to-real protocols answer different questions.  Our setting is target
offline adaptation: target interactions are counterfactual simulator labels
collected before deployment, and no target closed-loop reward is used for model
selection.

Mechanism-factored world models provide an intermediate design.  They restrict
prediction to stable parent sets and can adapt mechanism state separately,
consistent with the modularity motivation of causal transfer
\citep{bengio2020meta,feng2022factored,scholkopf2021causal}.  Our experiments
revealed a sharper
requirement for traffic control: accurate decomposition of absolute dynamics
does not guarantee accurate ordering of feasible phases.  The final method
therefore retains structural parent restrictions but changes the estimand from
next-state prediction to within-state action contrast.

We call the resulting method anchored Causal-Factored Cross-City Model
Transfer (CFCMT).  At a simulator snapshot, every feasible phase is executed
from the same restored state.  Candidate cost is expressed relative to one
source-selected pressure reference and normalized by the action range of that
matched group.  A low-capacity rigid anchor estimates candidate-versus-reference
cost from declared queue, occupancy, pressure and transition parents.  A
second model compares every unordered phase pair and is antisymmetrized by
construction.  The target score is
\begin{equation}
  S(s,a)=A(s,a)+\alpha(s)\{P(s,a)-A(s,a)\},
  \label{eq:intro-score}
\end{equation}
where the interpolation is selected out of fold by holding out one complete
target adaptation seed.  Constant interpolation permits aggressive correction
when it is stable; a local confidence-cap limits how far the pairwise model can
perturb the anchor when they disagree.

The information budget is central to the claim.  CFCMT is not evaluated as
zero-shot transfer in the final external experiment.  It uses target-network
files, online target state and counterfactual transition labels generated by an
uncalibrated target SUMO model.  Those labels constitute target offline
adaptation even though no target simulator parameter is fitted to field data.
The method is therefore calibration-light rather than target-data-free.
Adaptation seeds, the offline confirmation seed and closed-loop seeds are
disjoint, and the evaluation cache did not exist when the target model and
selector were frozen.

The empirical design has four levels.  First, 18 source networks from seven
city groups support leave-one-city-out development with an excluded evaluation
seed.  Second, two external city-derived networks, Los Angeles and Jinan, test
the frozen procedure using all valid target adaptation groups and a new
offline seed.  Third, 56 matched 3,600-s closed-loop rollouts compare CFCMT with
fixed time, simulator-only model predictive control, a rigid causal anchor, an
H2O$^+$-style dense residual and two classical pressure policies.  All policies
share target networks, feasible green phases, demand seeds, signal-transition
execution and metric populations.  Fourth, a target-label-only model already
frozen in the same baseline family is compared with CFCMT on 64 newly generated
seeds, four target scenarios and 512 rollouts.  This final stage directly tests
whether source-labelled counterfactuals improve closed-loop control after the
target adaptation budget is held fixed.  Official RESCO and LibSignal reinforcement
learning runs are retained only as native-protocol integration checks because
their state, reward, action timing and training budgets are not commensurate
with this experiment \citep{ault2021reinforcement,mei2023libsignal}.
The study consequently does not present native CrossLight, X-Light or GESA
scores as if they were paired baselines; a same-protocol reproduction of those
systems remains outside the evidence reported here.

This study makes four contributions.  First, it formulates cross-network
signal adaptation around matched action contrasts and proves cancellation of
action-common domain bias.  Second, it combines a parent-restricted anchor with
an antisymmetric all-action correction and derives exact ranking and bounded
perturbation properties.  Third, it implements an auditable seed-blocked
selection and full-refit protocol that prevents out-of-fold or evaluation
models from entering deployment.  Fourth, it separates structured-model gain
from source-data gain.  CFCMT improves the dense residual, but a fresh
target-label-only comparison reveals positive source contribution in Jinan and
negative transfer in Los Angeles, rejecting a general source-benefit claim.
Classical pressure control remains the closed-loop ceiling, and a separate
64-seed execution-layer challenge does not close that gap.  The contribution
is consequently an auditable action-ranking framework and a falsifiable account
of where cross-network transfer fails, not a claim that complexity or source
volume alone should replace a strong traffic-control rule.
